
% ============================================================
% SECTION 04
% INVARIANCE AND CHANGE
% Chapter 02: Invariance
% Book I: Ontology of Viability
% ============================================================

\section{Invariance and Change}
\label{sec:invariance-and-change}

% ------------------------------------------------------------
% 4.1. THE PROBLEM OF CHANGE IN THE ANALYSIS
% OF INVARIANCE
% ------------------------------------------------------------

The relationship between invariance and change constitutes
one of the central problems of the present chapter. The
relational definition of invariance adopted in this first
edition requires the specification of a domain of changes
or transformations in relation to which an existent, or
one of its determinate aspects, remains unchanged in a
particular respect.

Change is therefore directly relevant to the formulation
of invariance. However, their relationship must not be
understood through a simple opposition between two
mutually exclusive conditions.

If invariance were defined as the complete absence of
change, it would become indistinguishable from absolute
immutability. Such a definition would exclude the
possibility that an existent could undergo transformation
while some of its determinations remain unchanged.

Conversely, if every occurrence of change were understood
as the negation of all invariance, it would become
impossible to distinguish between the aspects of an
existent that vary and those that remain unchanged.

The analysis must therefore move beyond the general
opposition between permanence and change. It must
identify the respects in which change occurs and
the respects in which invariance may be asserted.

The purpose of this section is to clarify this relationship
without prematurely determining the ultimate ontological
ground of change, persistence, or invariance.

% ------------------------------------------------------------
% 4.2. CHANGE AS A DETERMINATE ALTERATION
% ------------------------------------------------------------

Change cannot be adequately understood as an
undifferentiated event in which an existent simply
becomes different.

A meaningful account of change requires the
identification of what changes, in what respect
it changes, and under which conditions the alteration
is considered.

An existent may undergo changes in position,
configuration, quantity, quality, relation, or
other determinations. These changes are not
necessarily equivalent, and they do not necessarily
affect every aspect of the existent in the same way.

For the purposes of this analysis, change is
considered as an alteration in one or more
determinations of an existent under specified
conditions.

This formulation is analytical. It does not
claim that every form of change can be reduced
to an alteration of properties, nor does it
assume that the distinction between an existent
and its determinations is an ontological
separation.

The distinction allows the analysis to examine
change without assuming that the existent as
a whole must be understood as either entirely
unchanged or entirely transformed.

% ------------------------------------------------------------
% 4.3. CHANGE DOES NOT NECESSARILY ELIMINATE
% INVARIANCE
% ------------------------------------------------------------

An existent may undergo change while some
determination remains unchanged in relation
to the same transformation.

Let \(A\) represent an existent and \(T\)
a specified transformation. Let \(F\) and
\(G\) represent two determinations considered
in relation to \(A\).

It is possible that:

\[
F(T(A)) \neq F(A)
\]

while:

\[
G(T(A)) = G(A)
\]

In this schematic representation, determination
\(F\) changes while determination \(G\) remains
unchanged under the specified transformation.

The expressions illustrate that change and
invariance may coexist when they concern
different determinations of the same existent.

This coexistence does not establish that
\(G\) is essential to the identity of \(A\),
nor does it demonstrate that the identity
of \(A\) necessarily persists through every
possible transformation.

It establishes only that a change in one
determination does not logically require
a change in every other determination.

The distinction is important because the
scope of a change must not be extended
beyond what the relevant analysis can
justify.

% ------------------------------------------------------------
% 4.4. THE DIFFERENCE BETWEEN LOCAL AND
% GENERAL CLAIMS OF CHANGE
% ------------------------------------------------------------

A claim about change may concern a particular
determination or the existent considered
more generally.

These are different levels of analysis.

A local claim identifies an alteration in
a specified determination. A general claim
concerns the extent to which the existent
as a whole has undergone transformation.

The first does not automatically establish
the second.

For example, a change in the spatial position
of an existent establishes a difference in
position. It does not establish that every
structural, relational, or other determination
has also changed.

Likewise, the preservation of a particular
relation does not establish that the existent
as a whole has remained unchanged.

The analysis must therefore distinguish
between:

\begin{enumerate}
    \item Change in a particular determination.
    \item Change in several determinations.
    \item Transformation of the overall
    configuration of an existent.
    \item A possible transformation affecting
    the identity of the existent itself.
\end{enumerate}

These distinctions identify different questions.
They do not establish a universal hierarchy
among the listed forms of change.

In particular, the fourth question cannot
be answered merely by accumulating observations
about changes in individual determinations.
It requires an account of the relationship
between those determinations and the identity
of the existent.

% ------------------------------------------------------------
% 4.5. INVARIANCE IS DOMAIN-DEPENDENT
% ------------------------------------------------------------

An invariance claim is restricted by the
domain of transformations in relation to
which it is formulated.

A determination may remain unchanged under
one set of transformations while changing
under another. Consequently, invariance
must not be interpreted as an unrestricted
claim that a determination remains unchanged
under every conceivable condition.

Let \(D_1\) and \(D_2\) represent two
specified domains of transformations.
A determination \(F\) may remain unchanged
under transformations belonging to \(D_1\),
while undergoing alteration under
transformations belonging to \(D_2\).

Schematically:

\[
F(T(A)) = F(A)
\qquad
\text{for } T \in D_1
\]

but there may be a transformation
\(T' \in D_2\) such that:

\[
F(T'(A)) \neq F(A)
\]

These expressions represent a restricted
invariance claim. They do not establish
that the domains are necessarily independent
or that every transformation can be assigned
to a formally defined mathematical set.

The important point is that invariance
and change must be examined relative to
specified domains.

Without such specification, the statement
that something is invariant remains
ambiguous, while the statement that
something changes may be too general
to support a precise ontological conclusion.

% ------------------------------------------------------------
% 4.6. CHANGE AND THE QUESTION OF
% CONTINUITY
% ------------------------------------------------------------

The relationship between change and
continuity must be distinguished from
the relationship between change and
invariance.

An existent may undergo a sequence of
changes, and the question may arise as
to whether it continues through those
changes.

However, the continuity of an existent
does not necessarily require that every
determination remain invariant.

Similarly, the invariance of a particular
determination does not by itself establish
the continuity of the existent as a whole.

These concepts address related but
different problems:

\begin{itemize}
    \item Change concerns alteration in
    one or more determinations.

    \item Invariance concerns the persistence
    of a specified determination relative
    to a defined domain of change.

    \item Continuity concerns whether and
    how an existent continues through
    different conditions or phases.
\end{itemize}

The relationship among these concepts
cannot be settled by treating any one
of them as a substitute for the others.

The ontological ground of continuity
and its relation to identity remain
open questions. They require a more
comprehensive account than the present
relational analysis of invariance can
provide.

% ------------------------------------------------------------
% 4.7. CHANGE DOES NOT NECESSARILY
% IMPLY LOSS OF IDENTITY
% ------------------------------------------------------------

The occurrence of change does not
automatically establish the loss of
identity.

The preceding chapter established
that being and identity are inseparable.
This commitment prevents the analysis
from treating existence as an initially
identity-neutral condition that may
simply lose identity whenever one of
its determinations changes.

At the same time, this foundational
commitment does not provide an immediate
criterion for deciding which transformations
are compatible with the persistence
of a particular existent and which,
if any, may involve a transformation
of identity.

The analysis must therefore distinguish
between the occurrence of change and
the ontological interpretation of that
change.

A change in position, quantity, configuration,
or relation may be identified without
having established what that alteration
means for the identity of the existent.

The question of identity through change
must remain open until the relationship
between identity, determination, and
the dynamics of being has been developed
more fully.

It would be premature to explain this
relationship by assuming an immutable
substrate, a permanent essence, or a
fixed collection of properties.

% ------------------------------------------------------------
% 4.8. CHANGE, OBSERVATION, AND
	EPISTEMIC LIMITATIONS
% ------------------------------------------------------------

The recognition of change frequently
depends upon the comparison of different
conditions or states.

An observer identifies an alteration
by distinguishing between what is
observed at one point or under one
condition and what is observed under
another.

This comparative activity belongs
to the epistemic dimension of the
analysis. It requires criteria for
identifying the relevant determinations
and distinguishing actual alteration
from differences in description,
measurement, or perspective.

However, the epistemic identification
of change must not be confused with
the ontological occurrence of change.

An existent may undergo an alteration
that remains unknown to a particular
observer. Likewise, an apparent change
may result from an inadequate comparison
or an error in the identification of
the relevant determination.

The same distinction applies to
invariance. The observer's inability
to identify a change does not establish
that no change has occurred. Nor does
the inability to identify an invariant
feature establish that no such feature
exists.

Knowledge and non-knowledge are
intertwined in the situated condition
of the knower. The analysis must
acknowledge this condition without
reducing the ontological status of
change or invariance to the limits
of observation.

% ------------------------------------------------------------
% 4.9. THE LIMITS OF THE RELATIONAL
% ACCOUNT OF CHANGE
% ------------------------------------------------------------

The relational definition of invariance
makes it possible to distinguish
between what changes and what remains
unchanged within a specified domain.

Nevertheless, it does not explain
why change occurs, what ultimately
constitutes transformation, or how
the identity-bearing existence of
an existent is related to its changing
determinations.

Several questions therefore remain
outside the scope of the present
section:

\begin{itemize}
    \item What is the ontological ground
    of change?

    \item Are all forms of change
    reducible to alterations in
    determinations, or are there
    other ontological forms of
    transformation?

    \item How should the relation
    between identity and change
    be understood?

    \item Under what conditions,
    if any, does a transformation
    involve a change in identity?

    \item Is invariance merely relational,
    or does it disclose a deeper
    dimension of ontological
    determination?

    \item How do the different realms
    of being participate in or
    condition the possibilities
    of change and invariance?
\end{itemize}

These questions are not answered
by assuming that change is always
external to an existent or that
invariance is always internal to it.
Such distinctions would require
additional ontological justification.

The present analysis is limited to
clarifying the conceptual relationship
between change and invariance.

% ------------------------------------------------------------
% 4.10. SYNTHESIS AND TRANSITION
% ------------------------------------------------------------

The relationship between invariance
and change cannot be adequately
described as a simple opposition
between permanence and alteration.

Change concerns the alteration
of one or more determinations under
specified conditions. Invariance
concerns the relation between an
existent, or one of its determinate
aspects, and a specified domain of
transformations, within which
something remains unchanged in
a particular respect.

These concepts may therefore coexist
within the analysis of the same
existent, provided that the relevant
determinations and respects are
properly distinguished.

The present section establishes
the following restrictions:

\begin{enumerate}
    \item Change in one determination
    does not necessarily imply
    change in every determination.

    \item Invariance in one respect
    does not establish the immutability
    of the existent as a whole.

    \item Change does not automatically
    imply the loss of identity.

    \item The persistence of an invariant
    determination does not by itself
    establish the continuity of the
    existent.

    \item Every claim of invariance
    must specify the domain of
    transformations in relation
    to which it is asserted.
\end{enumerate}

The analysis has therefore clarified
the conceptual conditions under
which change and invariance may
be considered together without
reducing either concept to the other.

The next section examines the
relationship between invariance
and continuity. This examination
is necessary because the persistence
of a determination under change
and the continuation of an existent
through different conditions are
related problems, but they cannot
be assumed to have the same
ontological meaning.

% ============================================================
% END OF SECTION 04
% ============================================================
